\subsection{Księga pierwsza -- definicje}
% Definicja 1.
\begin{euclid_definition}
    Punkt to jest to, co nie składa się z części.
\end{euclid_definition}

% Definicja 2.
\begin{euclid_definition}
    Linia jest długością bez szerokości.
\end{euclid_definition}

% Definicja 3.
\begin{euclid_definition}
    Końcami linii są punkty.
\end{euclid_definition}

% Definicja 4.
\begin{euclid_definition}
    Linia jest prosta, jeżeli położona jest między swoimi punktami w równym i jednostajnym kierunku.
\end{euclid_definition}

% Definicja 5.
\begin{euclid_definition}
    Powierzchnia jest to, co ma tylko długość i szerokość.
\end{euclid_definition}

% Definicja 6.
\begin{euclid_definition}
    Krawędzie powierzchni są liniami.
\end{euclid_definition}

% Definicja 7.
\begin{euclid_definition}
    Płaska powierzchnia albo płaszczyzna jest ta, na której biorąc gdziekolwiek dwa punkty linia prosta między tymi punktami cała leży na tej powierzchni.
\end{euclid_definition}

% Definicja 8.
\begin{euclid_definition}
    Kąt płaski to nachylenie dwóch linii na płaszczyźnie w miejscu, w którym jedna spotyka drugą i nie leżą w linii prostej.
\end{euclid_definition}

% Definicja 9.
\begin{euclid_definition}
    Kiedy linie są proste i tworzą kąt, wtedy kąt zwany jest prostoliniowym.
\end{euclid_definition}

% Definicja 10.
\begin{euclid_definition}
    Kiedy linia prosta padająca na drugą linie prostą, tworzy z nią kąty przyległe równe między sobą, to każdy z kątów równych nazywamy prostym, a padająca linia prostą nazywa się prostopadłą do tej linii, na którą pada.
\end{euclid_definition}

% Definicja 11.
\begin{euclid_definition}
    Kąt rozwarty jest większy od kąta prostego.
\end{euclid_definition}

% Definicja 12.
\begin{euclid_definition}
    Kąt ostry jest mniejszy od kąta prostego.
\end{euclid_definition}

% Definicja 13.
\begin{euclid_definition}
    Kresem albo granicą jest to, na czym się dana rzecz kończy.
\end{euclid_definition}

% Definicja 14.
\begin{euclid_definition}
    Figurą nazywamy to co jest ograniczone granicą lub granicami. 
\end{euclid_definition}

% Definicja 15.
\begin{euclid_definition}
    Koło jest figurą płaską zawarta linią zwaną okręgiem, do której wszystkie linie proste poprowadzone z jednego punktu wewnątrz figury położonego, są między sobą równe.
\end{euclid_definition}

% Definicja 16.
\begin{euclid_definition}
    I ten punkt nazywa się centrum lub środkiem koła. 
\end{euclid_definition}

% Definicja 17.
\begin{euclid_definition}
    Średnicą koła jest każda linia narysowana przez środek koła, przedłużona w dwóch kierunkach do jego obwodu, przepoławiająca go.
\end{euclid_definition}

% Definicja 18.
\begin{euclid_definition}
    Półokręgiem jest figura zawarta między średnicą i częscia okręgu odciętą tą średnicą. Środek półokregu jest też środkiem okręgu.
\end{euclid_definition}

% Definicja 19.
\begin{euclid_definition}
    Figury prostokreślne to figury ograniczone prostymi. Trójkąt to figura prostokreślna ograniczona trzema prostymi. Czworobok lub czworokąt to figura prostokreślna, która jest ograniczona czterema prostymi. Wielobok lub wielokąt to figura prostokreślna ograniczona więcej niż czterema prostymi.
\end{euclid_definition}

% Definicja 20.
\begin{euclid_definition}
    Trójkąt równoboczny to trójkąt, który ma trzy boki równe. Trójkąt równoramienny to trójkąt, który ma tylko dwa boki równe. Trójkąt różnoboczny to trójkąt, który ma trzy boki różne.
\end{euclid_definition}

% Definicja 21.
\begin{euclid_definition}
    Ponadto: trójkąt prostokątny to trójkąt, który na kąt prosty. Trójkąt rozwartokątny to trójkąt, który ma kąt rozwarty. Trójkąt ostrokątny to trójkąt, który ma trzy kąty ostre.
\end{euclid_definition}

% Definicja 22.
\begin{euclid_definition}
    Kwadrat jest to czworobok mający równe boki i równe kąty. Prostokąt jest to czworobok mający kąty proste, ale boki nierówne. Romb (kwadrat ukośny) jest to czworobok mający równe boki, ale nie mający kątów prostych. Równoległobok jest to czworobok mający boki przeciwległe równe, ale nie mający katów prostych. Wszystkie czworoboki inne niż wyżej wymienione nazywamy czworokątami.
\end{euclid_definition}

% Definicja 23.
\begin{euclid_definition}
    Linie równoległe, czyli mówiąc krócej równoległe są to proste, które leżą na tej samej płaszczyźnie i przedłużone z obu stron w nieskończoność, z żadnej strony nie przetną się.
\end{euclid_definition}